\documentclass[11pt, a4paper]{article}

% -- Paquetes ------------------------------------------------------------------
\usepackage[utf8]{inputenc}
\usepackage[T1]{fontenc}
\usepackage[spanish,es-noshorthands]{babel}
\usepackage{lmodern}
\usepackage[margin=2.5cm]{geometry}
\usepackage{amsmath, amssymb, amsthm}
\usepackage{graphicx}
\usepackage{tikz}
\usepackage{xcolor}
\usepackage{microtype}
\usepackage{parskip}
\usepackage[above,below]{placeins}
\usepackage{hyperref}

\hypersetup{colorlinks=true, linkcolor=black, citecolor=teal, urlcolor=teal}

\newcommand{\R}{\mathbb{R}}
\newcommand{\vOS}{\textsf{vOS}}
\newcommand{\nucleus}{\texttt{vsf-nucleus}}

% -- Título  -------------------------------------------------------------------
\title{%
  \textbf{Convergencia a Atractor Geométrico en un Kernel de Sistema Viable}\\[4pt]
  \large Resumen Ejecutivo
}

\author{Luis Pineda R.\thanks{Contacto: \texttt{luis@tallerpineda.cl}}}
\date{\textit{Informe Técnico -- Abril 2026; revisado Junio 2026}}

\begin{document}
\maketitle

% ??????????????????????????????????????????????????????????????????????????????
\section*{Qué Hicimos}

Construimos un kernel computacional (\nucleus{}) basado en el Modelo de Sistema Viable
(VSM) de Beer y probamos si un sistema gobernado por un conjunto reducido de axiomas
cibernéticos converge a un atractor estable en el espacio de estados. El estado del
kernel es un punto en $\R^{128}$, estructurado en bloques semánticos: identidad,
recursos, entropía, operaciones, canales de comunicación, memoria del atractor e
inteligencia regulatoria.

Ejecutamos doce experimentos sobre hardware físico (Intel Core i7-7600U, hipervisor
Proxmox) desde abril de 2026, midiendo convergencia, geometría de la cuenca,
resiliencia térmica, efectos del observador, ablación de la arquitectura y --- en el
experimento final --- el sistema completo ensamblado con todas las capas VSM activas.

% ??????????????????????????????????????????????????????????????????????????????
\section*{Resultados Principales}

\textbf{1. El atractor es un sumidero puntual estable, confirmado por tres vías
independientes.} Medimos el mayor exponente de Lyapunov directamente (método de
Benettin): $\lambda_\text{max} = -0{,}010050 = \ln(0{,}99)$ por tick, a precisión de
máquina --- negativo, de modo que las trayectorias cercanas se contraen en vez de
divergir. Una estimación Monte-Carlo sitúa el volumen de la cuenca en $\approx 1{,}0$
sobre la región alcanzable, y la homología persistente de la topología da
$\beta_0 = 1,\ \beta_1 = 0$ (un componente, sin agujero, sin ciclo límite). Dinámica,
volumen y forma coinciden: un único punto estable.

\textbf{2. Un certificado de viabilidad libre de coordenadas, $V = \min(g, c, t, m)$.}
Destilamos las mediciones en un funcional construido a partir de invariantes
intrínsecos (factibilidad, signo de contracción, firma topológica, margen de ruido).
En este kernel el eje ligante es $g = 0{,}102$ (la restricción de presupuesto de CPU),
y romper cualquier eje individual lleva $V \to 0$ por separado --- de modo que $V$ es
un vector discriminante, no un escalar cosmético siempre cercano a uno.

\textbf{3. El observador es constitutivo, no pasivo.} La observación estructurada baja
de forma medible la entropía del sistema en $0{,}31$ bits y comprime su dimensión
efectiva de $\approx 5{,}0$ a $3{,}9$ --- una instancia medida de cibernética de
segundo orden, en la que el observador es parte de la dinámica que estabiliza.

% ??????????????????????????????????????????????????????????????????????????????
\section*{Hallazgo Clave: Viabilidad en Dos Capas}

Los experimentos de ablación revelan una \textbf{estructura de viabilidad en dos
capas}:

\begin{figure}[h]
\centering
\begin{tikzpicture}[scale=1.1]
  \draw[thick, fill=blue!20] (0,3) rectangle (4,3.8);
  \node at (2,3.4) [text=white, font=\bfseries] {Capa 1: Arquitectura};
  \node at (2,2.5) {\small Bloques A--G construyen la cuenca};

  \draw[thick, fill=green!20] (0,1.5) rectangle (4,2.3);
  \node at (2,1.9) [text=white, font=\bfseries] {Capa 2: Gobernanza};
  \node at (2,0.9) {\small S5 aporta viabilidad condicional};

  \draw[dashed, thick] (0,0) .. controls (0.5,-0.3) and (3.5,-0.3) .. (4,0);
  \node at (2,-0.7) {\small Cuenca del atractor: piso $\cos \approx 0{,}875$ (peor caso)};
\end{tikzpicture}
\caption{La arquitectura construye la cuenca y su homeostasis; la gobernanza es la
diferencia entre la supervivencia y el colapso bajo estrés.}
\label{fig:dos_capas}
\end{figure}

\begin{description}
  \item[\textbf{Capa 1: Arquitectura}] La disposición de los bloques VSM determina
        físicamente qué regiones del espacio de estados son habitables. El sustrato
        desnudo (sin bloques) deriva fuera del objetivo y nunca se estabiliza; los
        bloques son lo que crea el sumidero puntual en $\R^{128}$.
  \item[\textbf{Capa 2: Gobernanza (S5)}] En reposo, la gobernanza es
        \emph{redundante} --- el sistema converge igual con o sin ella. Bajo estrés es
        decisiva. La gobernanza tiene dos brazos: uno \emph{reactivo} (la respuesta
        algedónica de seguridad que regula una carga térmica desbocada) y uno
        \emph{deliberativo} (el veto constitucional del pacto que puede detener la
        propia autonomía del sistema). Los experimentos ejercitan ambos.
\end{description}

Con ambas capas presentes el sistema mantiene la viabilidad en $\cos \approx 0{,}875$
(en escala $0$--$1$, con $0{,}70$ = umbral de viabilidad) incluso bajo estrés
impulsivo máximo --- un margen de $0{,}175$ nunca traspasado en ningún experimento.

% ??????????????????????????????????????????????????????????????????????????????
\section*{El Cierre: un Sistema Operacionalmente Cerrado}

El experimento final ejecutó el kernel completo con los cinco sistemas de Beer
(S1--S5) activos, bajo una sola perturbación. Cada nivel actuó dentro de su propio
ámbito cibernético. La observación decisiva atraviesa todo el modelo en un solo tick:
operaciones (S1) registró una caída, coordinación (S2) entregó una señal de política,
control (S3) detectó la degradación e \emph{intentó} una corrección, y política (S5)
la \emph{vetó} --- tras lo cual el sistema se autorrecuperó. Una prueba aparte
confirmó el veto constitucional de extremo a extremo: detiene el control autónomo bajo
degradación sostenida y lo restituye sólo ante una reanudación humana explícita y
registrada criptográficamente.

Éste es el sentido de la arquitectura: no una simulación del diagrama de Beer, sino un
sistema en el que cada flecha carga un evento medido.

% ??????????????????????????????????????????????????????????????????????????????
\section*{Resultados Empíricos}

\begin{tabular}{lrl}
\textbf{Propiedad} & \textbf{Medida} & \textbf{Estado} \\
\hline
Exponente de Lyapunov     & $\lambda_\text{max} = \ln(0{,}99) < 0$ & Medido, sumidero \\
Volumen de la cuenca      & $\approx 1{,}0$ (región alcanzable)    & Monte-Carlo \\
Topología                 & $\beta_0 = 1,\ \beta_1 = 0$            & Sin ciclo límite \\
Unicidad de la cuenca     & $\cos_\text{prom.\ pareado} = 0{,}9854$ & Confirmado \\
Viabilidad térmica        & viable a $68\,^\circ\mathrm{C}$        & Se mantiene \\
Piso de la cuenca         & $\cos \approx 0{,}875$                 & Margen $0{,}175$ \\
Caída de entropía         & $\Delta H = 0{,}31$ bits               & Constitutivo \\
Certificado de viabilidad & $V = \min(g,c,t,m),\ g = 0{,}102$      & Discrimina \\
Gobernanza de procesos    & 4 procesos del SO, $\tau = 1{,}0$      & cgroup~v2 vivo \\
Cierre operacional VSM    & S1--S5 causales en una corrida         & Demostrado \\
\end{tabular}

% ??????????????????????????????????????????????????????????????????????????????
\section*{Qué Significa Esto}

\begin{enumerate}
  \item \textbf{La viabilidad es geométrica.} La estabilidad del sistema emerge de la
        geometría de su espacio de estados, no de una imposición externa. Esto es lo
        que Beer entendía por un \emph{sistema viable}.

  \item \textbf{Las restricciones éticas entran como condiciones de factibilidad.} Los
        axiomas que prohíben el daño y la coerción no están añadidos por encima --- una
        violación del pacto colapsa la viabilidad ($V \to 0$) exactamente igual que
        agotar el presupuesto de CPU. La matemática es \emph{indiferente} a qué
        restricciones codifican ética; esa indiferencia es justamente lo que la hace un
        certificador general en vez de una teoría moral.

  \item \textbf{La gobernanza es resiliencia condicional.} Una buena capa de política
        no \emph{crea} la viabilidad --- eso lo hace la arquitectura --- pero es la
        diferencia entre la supervivencia airosa y el colapso cuando el sistema es
        desafiado.

  \item \textbf{La cuenca es profunda.} Bajo toda perturbación probada el sistema nunca
        se acercó al umbral de viabilidad ($0{,}70$); la aproximación más cercana fue
        $0{,}875$, un margen del $12{,}5\%$ sobre la falla.
\end{enumerate}

% ??????????????????????????????????????????????????????????????????????????????
\section*{Límites Honestos}

El trabajo corre sobre un único nodo de hardware; el estresor térmico, aunque eleva la
temperatura real del silicio, es una carga sintética; y la mayoría de las corridas son
de sistema cerrado. El límite final, y más nítido, es informativo: como el atractor
mantiene su medida de viabilidad invariante incluso ante una carga de CPU de más del
doble del número de núcleos, el canal \emph{anticipatorio} del sistema no puede leer
una alerta temprana de su propio estado regulado. La anticipación genuina requiere un
entorno externo --- variables cuya varianza sea independiente del propio atractor del
kernel. Acoplar el kernel a tal entorno es la dirección principal del trabajo en curso.

% ??????????????????????????????????????????????????????????????????????????????
\section*{Próximos Pasos}

\begin{enumerate}
  \item \textbf{Acoplamiento a un entorno externo.} Conectar el kernel a unidades
        operativas genuinamente externas (el caso motivante: una red distribuida de
        sensado climático) cuyas señales no estén supeditadas al atractor del kernel
        --- la prueba que el canal anticipatorio realmente necesita.
  \item \textbf{Coordinación inter-sistema vía un sustrato formal.} Una línea de
        trabajo complementaria estudia si dos modelos heterogéneos pueden coordinarse
        como capas VSM distintas a través de un lenguaje formal compartido (base de un
        segundo artículo).
  \item \textbf{\vOS{} sobre metal desnudo.} Desplegar el kernel contra cargas de
        trabajo reales (bases de datos, inferencia) para medir la geometría del
        atractor bajo perfiles de carga diversos.
\end{enumerate}

% ??????????????????????????????????????????????????????????????????????????????
\section*{Para la Comunidad Beer}

Este trabajo vuelve a los primeros principios: \emph{¿qué significa que un sistema sea
viable?} Damos una respuesta computacional anclada en axiomas cibernéticos, con
instrumentos --- el algoritmo de Benettin, la homología persistente, la atestación
criptográfica --- que no existían cuando Beer escribió. La distinción
arquitectura--gobernanza, y la demostración de que los cinco sistemas cierran un bucle
operacional sobre hardware real, sugieren que los bloques del VSM no son metáforas: son
requisitos de viabilidad con geometría medible.

El código es de fuente abierta. El artículo completo, los datos y las bitácoras
experimentales están disponibles a pedido.

\end{document}
